% ===================================================================
\section{Recorridos guiados}

Tres recorridos cortos que muestran los resultados principales del proyecto.

\subsection{De la estabilidad a la cascada (10 minutos)}
\begin{enumerate}
\item En \emph{Inicio}, presione \boton{Volver a la red por defecto}.
\item En \emph{Diagnóstico}, deje $\alpha=1$ y mueva $\beta$ hasta $R\approx0.8$. El
semáforo está en verde.
\item Vaya a \emph{Simulador de choques}. Choque la entidad que más paga con un retiro de
3 ciclos y colchón 0.5. Anote el porcentaje en impago.
\item Vuelva a \emph{Diagnóstico}, suba $\beta$ hasta $R\approx1.2$ (semáforo rojo) y repita
el choque. La cascada es mucho mayor.
\item Suba el colchón a 1.2 y repita. La cascada se reduce: el colchón es la defensa de
cada entidad.
\end{enumerate}

\subsection{La ventana de contagio (5 minutos)}
\begin{enumerate}
\item En \emph{Explorar escenarios}, deje todos los filtros abiertos.
\item Observe la gráfica \emph{Ventana de contagio}: la fracción en impago es máxima con
grado medio intermedio.
\item Filtre solo choques de liquidez y luego solo operativos, y compare.
\end{enumerate}

\subsection{La sincronía como riesgo (10 minutos)}
\begin{enumerate}
\item En \emph{Cartera multicapa}, simule la cartera por defecto en el mundo con red y
coloree por coherencia $\hat r$. Los formales aparecen verdes; los informales, amarillos o
rojos.
\item Revise los diagramas de caja de \emph{¿Quién está desincronizado?}.
\item Simule la quiebra del mayor empleador y compare con el mundo nulo.
\item En \emph{Riesgo crediticio híbrido}, lea las tres métricas superiores: la dinámica
mejora a la red estática en el mundo con red y no mejora nada en el nulo.
\item Abra el cliente con mayor riesgo oculto y lea sus barras SHAP.
\end{enumerate}

% ===================================================================
\section{Solución de problemas}

\begin{longtable}{p{0.36\textwidth}p{0.58\textwidth}}
\toprule
\textbf{Síntoma} & \textbf{Causa y solución} \\
\midrule
\endhead
\emph{Red demasiado grande para simulación en vivo} & La red tiene más de 300 entidades.
Redúzcala en \emph{Construir red} o use el \emph{Mapa de riesgo}. \\
\emph{El grado medio debe ser menor que n - 1} & El grado pedido es imposible para ese
tamaño. Baje el grado o suba el número de entidades. \\
\emph{No hay dataset} (Explorar escenarios) & Ejecute \texttt{uv run pulso sweep
configs/sweep.yaml}. \\
\emph{No hay modelos entrenados} (Mapa de riesgo) & Ejecute \texttt{uv run pulso train
configs/train.yaml}. \\
\emph{No hay resultados del experimento} (páginas 8 y 9) & Ejecute \texttt{uv run pulso
credit configs/credit.yaml}. \\
La página no cambia al mover un control & Si el control está dentro de un formulario,
presione su botón azul. Espere a que desaparezca el indicador de ejecución. \\
Los valores volvieron a los iniciales & Se recargó la pestaña del navegador. El estado
dura solo mientras la pestaña esté abierta. \\
\emph{Port 8501 is already in use} & Otra aplicación usa el puerto. Agregue
\texttt{--server.port=8502}. \\
\emph{ModuleNotFoundError: No module named 'pulso'} (macOS) & macOS oculta un archivo de
la instalación dentro de \texttt{\textasciitilde/Documents}. Ejecute
\texttt{chflags -R nohidden .venv} o anteponga \texttt{PYTHONPATH=src} al comando. \\
Una simulación tarda mucho & Redes grandes o tiempos largos. La primera vez cada
cálculo tarda más; después queda guardado en memoria. \\
\bottomrule
\end{longtable}

% ===================================================================
\section{Glosario rápido}

\begin{longtable}{p{0.24\textwidth}p{0.7\textwidth}}
\toprule
\textbf{Término} & \textbf{Significado} \\
\midrule
\endhead
$\alpha$ & Velocidad con que una entidad corrige su déficit. \\
$\beta$ & Intensidad del contagio de déficit entre pagador y receptor. \\
$\gamma$ & Déficit que genera el ingreso desfasado. \\
$\Delta$ & Dispersión de los ritmos naturales de pago. \\
$\lambda$ & Déficit que genera perder un pagador. \\
$\lambda_2$ & Conectividad algebraica de la red. \\
$\sigma$ & Estabilidad: negativa, la red se recupera; positiva, amplifica. \\
$\tau_\text{rec}$ & Tiempo de recuperación tras un choque. \\
$D_0$ & Intensidad de la redistribución de liquidez. \\
$K$ & Intensidad con que se ajustan los ritmos de pago. \\
$R$ & Índice de amplificación; menor que 1, la red amortigua. \\
$r$, $\hat r$ & Sincronía real y estimada desde transacciones. \\
AUC & Capacidad de un modelo para ordenar a los que caen por encima de los que no. \\
Colapso global & Impago del 30~\% o más de la red. \\
DebtRank & Medida del daño sistémico de un choque. \\
GCN & Red neuronal que trabaja sobre grafos. \\
Mundo nulo & Control en que la red no transmite déficit. \\
PD & Probabilidad de impago a 12 meses. \\
SHAP & Método que explica cuánto aporta cada variable a una predicción. \\
SMMLV & Salario mínimo mensual legal vigente. \\
XGBoost & Modelo de aprendizaje automático basado en árboles de decisión. \\
\bottomrule
\end{longtable}
